\documentclass[../../../include/open-logic-section]{subfiles}

\begin{document}
	
\olfileid[pt]{sth}{replacement}{finiteaxiomatizability}
\olsection{Apêndice: Axiomatizabilidade finita}
Terminamos este capítulo extraindo alguns resultados da Substituição. O primeiro resultado é devido a \citet{Montague1961}; note que não é uma prova \emph{dentro} de $\ZF$, mas uma prova \emph{sobre} $\ZF$:
\begin{thm}\ollabel{zfnotfinitely}
	$\ZF$ não é finitamente axiomatizável. Mais geralmente: se $\Th{T}$ é finita e $\Th{T} \Proves \ZF$, então $\Th{T}$ é inconsistente.
	
	(Aqui nos restringimos tacitamente a sentenças de primeira ordem cujo único primitivo não lógico é $\in$, e escrevemos $\Th{T} \Proves \ZF$ para indicar que $\Th{T} \Proves \phi$ para todo o $\phi \in \ZF$.)
\end{thm}
\begin{proof}
	Fixemos $\Th{T}$ finita tal que $\Th{T} \Proves \ZF$. Logo $\Th{T}$ prova a Reflexão, O.s.\@, \olref[sth][replacement][ref]{reflectionschema}. Como $\Th{T}$ é finita, podemos reescrevê-la como uma única conjunção, $\theta$. Refletindo com esta fórmula, $\Th{T} \Proves \exists \beta(\theta \liff \theta^{V_\beta})$. Como trivialmente $\Th{T} \Proves \theta$, obtemos $\Th{T} \Proves \exists \beta\ \theta^{V_\beta}$. 
	
	Ora, seja $\psi(X)$ uma abreviatura de:
	\[
		\theta^X \land X\text{ é transitivo} \land (\forall Y \in X)(Y\text{ é transitivo}\lif \lnot \theta^{Y})
	\]
	isto diz, em traços largos: $X$ é um modelo transitivo de $\theta$, e $\in$-minimal a esse respeito. Ora, recordando que $\Th{T} \Proves \exists \beta\ \theta^{V_\beta}$, por fatos básicos sobre postos dentro de $\ZF$ e portanto dentro de $\Th{T}$, temos:
	\begin{equation}
		\Th{T} \Proves \exists M \psi(M). \tag{*}\label{Mpsi}
	\end{equation}
	Usando o primeiro conjunto de $\psi(X)$, sempre que $\Th{T} \Proves \sigma$, temos $\Th{T} \Proves \forall X(\psi(X) \lif \sigma^X)$. Logo, por \eqref{Mpsi}:
	\begin{align*}
		\Th{T} &\Proves \forall X(\psi(X) \lif (\exists N \psi(N))^X)\\
	\intertext{Usando isto, e \eqref{Mpsi} de novo:}
		\Th{T} &\Proves \exists M(\psi(M) \land (\exists N \psi(N))^M)
	\intertext{Em particular, então:}
		\Th{T} &\Proves \exists M(\psi(M) \land (\exists N \in M)((N\text{ é transitivo})^N \land (\theta^N)^M))
	\intertext{Logo, por raciocínio elementar sobre transitividade:}
		\Th{T} &\Proves \exists M(\psi(M) \land (\exists N \in M)(N\text{ é transitivo} \land \theta^N))
	\end{align*} 
	Logo $\Th{T}$ é inconsistente.\footnote{Este ``raciocínio elementar'' envolve provar certos ``fatos de absolutidade'' para conjuntos transitivos.}
\end{proof}

Eis um resultado semelhante, assinalado por \citet[223]{Potter2004}:

\begin{prop}\ollabel{finiteextensionofZ}
	Seja $\Th{T}$ uma extensão de $\Z$ com um número finito de novos axiomas. Se $\Th{T} \Proves \ZF$, então $\Th{T}$ é inconsistente. (Aqui usamos as mesmas restrições tácitas que para \olref{zfnotfinitely}.)
\end{prop}
\begin{proof}
	Usemos $\theta$ para a conjunção de todos os axiomas de $\Th{T}$ \emph{exceto} as (infinitas) instâncias de Separação. Definindo $\psi$ a partir de $\theta$ como no \olref{zfnotfinitely}, podemos mostrar que $\Th{T} \Proves \exists M \psi(M)$. 
	
	Como no \olref{zfnotfinitely}, podemos estabelecer o esquema de que, sempre que $\Th{T} \Proves \sigma$, temos $\Th{T} \Proves \forall X(\psi(X) \lif \sigma^X)$. Terminamos a prova exatamente como no \olref{zfnotfinitely}.
	
	No entanto, estabelecer o esquema envolve um pouco mais de trabalho do que no \olref{zfnotfinitely}. Afinal, as instâncias de Separação estão em $\Th{T}$, mas não são conjuntos de $\theta$. Contudo, podemos ultrapassar este obstáculo provando que $\Th{T} \Proves \forall X(X\text{ é transitivo} \lif \sigma^X)$, para toda a instância de Separação $\sigma$. Deixamos isto ao leitor. 
\end{proof}
\begin{prob}
	Mostre que, para toda a instância de Separação $\sigma$, tem-se: $\Z \Proves \forall X(X\text{ é transitivo} \lif \sigma^X)$. (Usamos este esquema na \olref[sth][replacement][finiteaxiomatizability]{finiteextensionofZ}.)
\end{prob}
\begin{prob}
	Mostre que, para todo o $\phi \in \Z$, tem-se $\ZF \Proves \phi^{V_{\omega+\omega}}$.
\end{prob}
\begin{prob}
	Confirme os restantes resultados esquemáticos invocados nas provas de \olref[sth][replacement][finiteaxiomatizability]{zfnotfinitely} e  \olref[sth][replacement][finiteaxiomatizability]{finiteextensionofZ}.
\end{prob}

Como observamos na \olref[sth][replacement][absinf]{sec}, isto mostra que a Substituição é estritamente mais forte que
\olref[ordinals][ordtype]{thmOrdinalRepresentation}. Ou, um pouco mais precisamente: se $\Z$ + ``toda a boa ordenação é isomorfa a um ordinal único'' for consistente, então não prova alguma instância da Substituição.


	%	By assumption, $\psi^{V_\alpha}$ has the form: 
	%	\[
	%		(\forall A \in V_\alpha)(\exists S \in V_\alpha)(\forall x \in V_\alpha)(x \in S \liff (\phi^{V_\alpha}(x) \land x \in A))
	%	\]
	%	To establish this holds, fix $A \in V_\alpha$. Using Separation, obtain: 
	%	\[
	%		S = \Setabs{x \in A}{\phi^{V_\alpha}(x)}
	%	\]
	%	Now $S \in V_\alpha$, since $S \subseteq A \in V_\alpha$, and clearly $(\forall x \in V_\alpha)(x \in S \liff (\phi^{V_\alpha}(x) \land x \in A)$.

\end{document}
